\section{Conclusion}
\label{sec:conclusion}

Short-Burst is a meta-algorithm: it wraps any ergodic Markov chain in a
burst-and-select loop to drive compactness optimization.
The chain type is a parameter.

We have implemented two new variants---\sbf{} and \sbms{}---by substituting
ForestRecom and MergeSplit for the standard \recom{} inner chain.
The substitution delivers approximate distributional correctness at the
burst endpoints (from the ForestRecom and MergeSplit MH corrections) while
preserving the compactness-optimization structure of the Short-Burst wrapper.

The empirical comparison on NC, WI, and TX reveals a clean tradeoff:
at equal compute budget, SB-Standard achieves the best final edge-cut because
its uncorrected \recom{} inner chain has higher acceptance rate ($\approx 75\%$)
and thus visits more distinct plans per burst;
\sbf{} achieves somewhat higher edge-cut but produces approximately correct
endpoint distributions, making it the preferred choice for ensemble
characterisation near the optimum;
\sbms{} sits between the two variants on both dimensions.

The per-burst acceptance rate diagnostic $\alpha_i$, recorded in the audit
manifest for all three variants, provides both a reproducibility integrity
check and an empirical characterisation of the effective burst length.
It is a lightweight addition to the manifest that costs no extra compute
and provides meaningful diagnostic information.

\paragraph{Implementation.}
All three variants are available in \texttt{redist-cli}:
\begin{verbatim}
--search short-burst            # G.6: standard ReCom inner chain
--search short-burst-forest     # G.12: ForestRecom inner chain
--search short-burst-merge-split # G.12: MergeSplit inner chain
\end{verbatim}

\paragraph{Composability.}
The chain-type parameter exposes the composable structure of Short-Burst
and suggests a general principle: any future chain type that operates on
the redistricting plan space---subject only to ergodicity and the ability
to generate forward-RNG and reverse-RNG seeds---can be plugged into the
burst wrapper with no modification to the burst-selection logic.
This architectural decision, made in G.6 and extended here, provides a
stable foundation for future ensemble method development in the G series.
